\documentclass[10pt,a4paper]{amsart}
\usepackage[margin=19mm]{geometry}
\usepackage{amsmath,amssymb,mathtools}
\usepackage[scr=rsfs,cal=euler]{mathalfa}
\usepackage{xfrac}
\usepackage[unicode,hidelinks,bookmarks=false]{hyperref}
\newcommand{\ie}{i.e.\ }
\newcommand{\cf}{cf.\ }
\newcommand{\resp}{respectively\ }
\newcommand{\Subsection}{\subsection}
\providecommand{\nobreakdash}{\nobreak\hbox{-}}
\setlength{\emergencystretch}{2em}
\multlinegap0pt
\usepackage{bm}
\usepackage{bbm}
\usepackage{dsfont}
\usepackage{xspace}
\usepackage{cancel}
\usepackage{mathtools}
\usepackage{amsthm}
\makeatletter
\@ifundefined{theorem}{\newtheorem{theorem}{Theorem}}{}
\makeatother
\newcommand{\D}{\mathbb{D}}
\newcommand{\N}{\mathbb{N}}
\newcommand{\T}{\mathbb{T}}
\title[Interpolating Sequences For Dual Uniform Algebras]{Interpolating Sequences For Dual Uniform Algebras}
\author{Mario P. Maletzki}
\date{Source: arXiv:2607.23796v1 (CC BY 4.0)}
\begin{document}
\maketitle

\noindent\emph{Source and licence.} Interpolating Sequences For Dual Uniform Algebras. Version arXiv:2607.23796v1 (CC BY 4.0), distributed under the Creative Commons Attribution 4.0 International licence, as recorded in the arXiv licence metadata for this exact version and shown on the abstract page https://arxiv.org/abs/2607.23796v1. Full rights notice: licenses/arxiv-2607.23796-CC-BY-4.0.txt.\par\smallskip

\noindent\textit{Editorial scope.} Counted unit: the Theorem whose source label is \texttt{None} in the article (source0.tex lines 512--518, source label \texttt{None}), delivered together with the article's own complete original proof of it (source0.tex lines 520--550, one \texttt{proof} environment of 31 lines). No mathematics is written, repaired or re-derived: statement and proof are the article's own text line for line. Required context retained uncounted: none. Every definition, display and label of those lines is retained unchanged. Omitted from this package and not redistributed: 1--511, 519--519, 551--926. Nothing inside a retained excerpt is omitted or altered: every retained body is delivered exactly as the article's lines are written, so no omission receipt of any kind accompanies this package. Citations: the retained excerpts carry no citation command at all, so no bibliography entry of the article is transcribed and no reference is redistributed. Reference adaptation: none was needed and none was made; every reference inside the delivered unit resolves inside it, so no unresolved reference or citation is produced. Printed slips kept literal: no printed word, symbol, hypothesis or number of the counted statement or of its proof was altered. Layout and class: the article's own class is \texttt{elsarticle} with packages including amssymb, mathrsfs, hyperref, amsmath, amsthm; the wrapper is the lane's pinned \texttt{amsart} editorial wrapper at 10pt on A4, which restates the theorem-like environments the retained text needs and adds the article's own macro definitions that the retained text needs (\texttt{\textbackslash D}, \texttt{\textbackslash N}, \texttt{\textbackslash T}), with extra wrapper packages (bm, bbm, dsfont, xspace, cancel, mathtools). The delivered numbering is the unit's own. Assets: no retained span contains a figure environment, a graphics-inclusion command, a PDF-page inclusion, a table or a TikZ picture; no image file is copied into this package, the package declares no asset dependency and no image file is redistributed with it. Licence: version arXiv:2607.23796v1 of this article is distributed under the Creative Commons Attribution 4.0 International licence, as recorded in the arXiv licence metadata for this exact version and shown on the abstract page https://arxiv.org/abs/2607.23796v1. A full rights notice is delivered at licenses/arxiv-2607.23796-CC-BY-4.0.txt. Characters: every retained excerpt is pure ASCII, so the delivered engine, XeLaTeX with its default Latin Modern text font, reports no missing character.
\begin{theorem}
    A sequence $(z_n)_n$ in $\D^N$ satisfies 
    \[
    \delta:=\inf_n \prod_{k\neq n}\rho_{\D^N}(z_k,z_n)>0
    \]
    if and only if $(z_n)_n$ is an interpolating for $H^\infty(\D^N)$ such that $\sum_{n=1}^\infty(1-\lVert z_n\rVert)<\infty$.
\end{theorem}

\begin{proof}
For each positive integer $n$ we consider the partition of $\N\setminus\{n\}$ into the sets $s_n(1), s_n(2), \dots, s_n(N)$ defined by letting $k$ be in $s_n(1)$ if $\rho_{\D^N}(z_n,z_k)=\rho_{\D}(z_n^1,z_k^1)$, and for $2\leq i\leq N$, letting $k$ be in $s_n(i)$ if $k\notin s_n(j)$ for $j<i$ and $\rho_{\D^N}(z_n,z_k)=\rho_{\D}(z_n^i,z_k^i)$. Using these partitions it is then easily seen that $\prod_{k\neq n}\rho_{\D^N}(z_k,z_n)>0$ for some $n$ if and only if $\sum_{n=1}^\infty(1-\lVert z_n\rVert)<\infty$. 

Now, if $(z_n)_n$ is interpolating, then by the Open Mapping theorem there are functions $(f_n)_n\subset H^\infty(\D^N)$ and a constant $M>0$ such that $f_n(z_k)=\delta_{n,k}$ for every pair of positive integers $k,n$ while $\lVert f_n \rVert<M$ for every $n$. Fixing a positive integer $m>1$, we then define for each $n$ the functions 
\[
\phi_n^m(z^1,\dots,z^N):=f_n(z^1,\dots,z^N)\prod_{i=1}^N\left(\prod_{\substack{k=1\\k\in s_n(i)}}^m \frac{1-\overline{z_k^i}z^i}{z^i-z_k^i} \right).
\] 
Each $\phi_n^m$ is holomorphic on $\D^N\setminus E_n$, where 
\[
E_n:=\bigcup_{i=1}^N\left(\bigcup_{\substack{k=1\\k\in s_n(i)}}^m\left(\D^{i-1}\times\{z_k^i\}\times \D^{N-i}\right)\right),
\] 
and we claim there is a holomorphic extension of $\phi_n^m$ to $\D^N$ that is bounded by $M$. In fact, if we consider $z_k^i$ for some $1\leq i\leq N$ and $k\in s_n(i)$ with $1\leq k\leq m$, then for every $(w^1,\dots,w^{i-1},z_k^i,w^{i+1},\dots,w^N)$ in $E_n$ such that no $w^j$ belongs to the set $\{z_k^j : k\in s_n(j), 1\leq k\leq m\}$, we have that $\varphi_{n,i}^m(z):=\phi_n^m(w^1,\dots,w^{i-1},z,w^{i+1},\dots,w^N)$ has removable singularities in $\{z_k^i : k\in s_n(i), 1\leq k\leq m\}$. Denoting also by $\varphi_{n,i}^m$ its holomorphic extension to $\D$, we define $\widetilde{\phi_n^m}=\phi_n^m$ outside $E_n$ and $\widetilde{\phi_n^m}(w^1,\dots,w^{i-1},z_k^i,w^{i+1},\dots,w^N):=\varphi_{n,i}^m(z_k^i)$. If we now consider a point $(w^1,\dots,w^{i-1},z_k^i,w^{i+1},\dots,w^N)$ such that there is exactly one $w^j$ belonging to $\{z_k^j : k\in s_n(j), 1\leq k\leq m\}$, then the function that maps 
\[
z\mapsto\widetilde{\phi_n^m}(w^1,\dots,w^{j-1},z,w^{j+1},\dots,w^{i-1},z_k^i,w^{i+1},\dots,w^N)
\] 
also has removable singularities in $\{z_k^j : k\in s_n(j), 1\leq k\leq m\}$, and thus we may define $\widetilde{\phi_n^m}$ in these points as the values of its holomorphic extension. Repeating the argument $N-2$ more times, we will have recursively extended the function $\phi_n^m$ to a holomorphic function $\widetilde{\phi_n^m}$ on $\D^N$. Moreover, by the Maximum Modulus Principle for $\D^N$ we have that $\lVert\widetilde{\phi_n^m}\rVert=\lim_{r\to 1}\max\{|\widetilde{\phi_n^m}(z^1,\dots,z^N)|: (z^i)_{i=1}^N\subset r\T\}$, so in particular it is $\lVert\widetilde{\phi_n^m}\rVert=\lVert f_n\rVert\leq M$. 
    
The condition $\lVert\widetilde{\phi_n^m}\rVert\leq M$ guarantees that 
\[
|f_n(z^1,\dots,z^N)|\leq M \left(\prod_{i=1}^N\left(\prod_{\substack{k=1\\k\in s_n(i)}}^m \left|\frac{z^i-z_k^i}{1-\overline{z_k^i}z^i}\right| \right)\right)\quad \text{for }\ (z^1,\dots, z^N)\in\D^N,
\] 
and in particular 
\[
1=|f_n(z_n^1,\dots,z_n^N)|\leq M \left(\prod_{i=1}^N\left(\prod_{\substack{k=1\\k\in s_n(i)}}^m \left|\frac{z_n^i-z_k^i}{1-\overline{z_k^i}z_n^i}\right| \right)\right).
\] 

Taking the limit as $m \to \infty$, the last product converges because of the hypothesis $\sum_{n=1}^\infty(1-\lVert z_n\rVert)<\infty$, and we conclude that 
\[
\prod_{\substack{k=1\\k\neq n}}^\infty \rho_{\D^N}(z_n,z_k)=\prod_{i=1}^N\left(\prod_{\substack{k=1\\k\in s_n(i)}}^\infty \left|\frac{z_n^i-z_k^i}{1-\overline{z_k^i}z_n^i}\right| \right)\geq \frac{1}{M}>0\quad \text{for every } n.
\]
\end{proof}

\end{document}
